\section{\uppercase{Introduction}}
\label{sec:introduction}
\noindent
During online handwriting acquisition, a digitizing tablet records the pen trajectory together with its speed, pressure and inclination. Encoding these measurements as color adds information about how a pen trace was produced while keeping its recognizable shape, so that standard image descriptors and classifiers can use the dynamics directly.

We aim to distinguish participants with Parkinson's disease (PD) from healthy participants (H) using PaHaW, a public annotated dataset containing online handwriting and drawing tasks from 75 participants \cite{Drotar2016}. Handwriting exposes aspects of fine motor behavior, and the relation between pen motion and diagnosis depends on the task performed, device used and participant characteristics \cite{Impedovo2019}.

We propose \emph{pseudodynamic minimaps}, a new image representation of online handwriting task executions, in which the planar coordinates place each point and color encodes the dynamics of the stroke passing through it. The term minimap follows work that renders source code compactly as a visual signature \cite{Gebara2025}.

Handwriting classification has often used explicitly designed kinematic and pressure features, as in the study that introduced PaHaW \cite{Drotar2016}. Learned features provide another route, either from pen dynamics \cite{Pereira2016} or from visual attributes of reconstructed writing \cite{Moetesum2019}, and dynamic information has also been carried by handwriting images themselves \cite{Diaz2019,Cilia2021SyntheticImages}.

Rendering a temporal signal as an image is an established route to reusing visual descriptors and convolutional models: recurrence plots expose the recurrence structure of a dynamical system as a two-dimensional pattern \cite{Eckmann1987}, Gramian angular fields and Markov transition fields encode a univariate series as a matrix \cite{Wang2015}, and convolutional networks trained on such renderings are competitive on time-series benchmarks \cite{Hatami2018}.

Pseudodynamic minimaps differ from these transforms and from the handwriting images above in three respects that can be checked against any of them. The image coordinates are the physical positions the pen visited, rather than a time axis or a synthetic layout, so each signal value stays anchored to the stroke that produced it: the closest handwriting encodings add the dynamics to a static trace as extra marks and pen-up strokes \cite{Diaz2019}, or arrange the signals in an image whose geometry is not the writing \cite{Cilia2021SyntheticImages}. Each signal is normalized within its own recording by percentiles, so color describes the writer's relative variation instead of a cohort-wide scale. Azimuth is centered on its circular mean, so that the arbitrary zero-degree boundary does not separate similar orientations. What we contribute is this formulation and a controlled account of what its performance depends on; the methods above are not re-implemented here, and the comparisons are with the static trace, kinematic attributes and pretrained features under one protocol.

This initial study evaluates the representation in three experiments that share participant-separated partitions: single signals and RGB triplets added to the static trace, a training-only selection of the encoding with controls that displace the signals along each stroke, and a comparison with kinematic attributes and pretrained features on the repeated \textit{l}, \textit{le} and \textit{les} tasks examined in earlier work on PaHaW \cite{Casademunt2023}.

The results indicate that the classifiers make use of the encoded dynamics: training-only selection preferred a dynamic encoding over the static trace in almost every training set, displacing the signals along the stroke lowered performance, and on the three repeated-pattern tasks the minimap pipeline gave the best observed participant-level decisions among the four pipelines. A cohort of 75 participants leaves the size of these gains uncertain; the study establishes how pen dynamics can be carried into an image and where they help, rather than delivering a clinical tool.
